\chapter{Intersection theory and b-divisors}
\index{intersection theory!b-divisor@$b$-divisor}\index{b-divisor@$b$-divisor!Cartier and Weil}
\label{app:intersection-bdivisors}

This appendix records the intersection-theoretic facts used on birational
models.
Its central firewall is categorical:
Cartier b-divisors are eventually pullback data, whereas Weil b-divisors are
compatible pushforward data.
Pullback and intersection are therefore automatic for Cartier objects and
are not automatic for Weil objects.

The first use is the model intersection construction in \cref{part:foundations}.
\Cref{chap:bj-homogeneity-bdivisors} uses the coefficient functions
of b-divisors, \Cref{chap:bj-cartier-realization} uses Cartier
realizations, and \Cref{chap:compare-ma} compares their valid
intersection products with analytic Monge--Amp\`ere products.
The Segre-degree example in \cref{sec:appF-Segre-ample-boundary}
explains why numerical detection of ideal singularities can require
an ample polarization even when a big and nef one is available.

The Chern-class discussion in \cref{sec:appF-Chern-volume-bridge}
connects ordinary projective intersections to the normalized
mixed zeta numerator of \cref{sec:appC-mixed-Segre-zeta}.

\Cref{sec:appF-total-Chern-volume,sec:appF-Richardson-volume}
then develop the same mechanism for total Chern intersections
and tableau coefficients. The volume operators of
\cref{sec:appD-volume-operators} keep track of the required factorials.

\section{Cycles, Cartier divisors, and numerical classes}
\label{sec:appF-intersection-basics}

Let \(X\) be a proper integral variety of dimension \(n\) over a field.
Write \(Z_k(X)\) for \(k\)-cycles with real coefficients and
\(A_k(X)_{\mathbf R}\) for cycles modulo rational equivalence.
An \(\mathbf R\)-Cartier divisor is a finite real linear combination of
Cartier divisors.

\begin{definition}[Numerical equivalence]
\label{def:appF-numerical}
Two \(\mathbf R\)-Cartier divisors \(D,D'\) are numerically equivalent if
\[
                     (D-D')\cdot C=0
\]
for every integral curve \(C\subset X\).
The quotient is \(N^1(X)\).
The real vector space of one-cycles modulo numerical equivalence is
\(N_1(X)\), and intersection gives a pairing
\[
                           N^1(X)\times N_1(X)\longrightarrow\mathbf R.
\tag{F.1}\label{eq:appF-N1-pairing}
\]
\end{definition}

A class \(\alpha\in N^1(X)\) is nef if
\(\alpha\cdot C\geq0\) for every integral curve.
It is ample only when represented by an ample \(\mathbf R\)-divisor in the
usual open cone.
Nefness is numerical; effectivity is not.

\begin{proposition}\label{prop:appF-mixed-cartier}
For \(\mathbf R\)-Cartier divisors \(D_1,\ldots,D_n\) on \(X\), the degree
\[
                              (D_1\cdots D_n)
\tag{F.2}\label{eq:appF-mixed-intersection}
\]
is symmetric and multilinear and depends only on their numerical classes.
If all \(D_i\) are nef, the number is nonnegative.
\end{proposition}

\begin{proof}
For integral Cartier divisors, intersect successively with the fundamental
cycle in the Chow group.
Commutativity of Cartier divisor operators gives symmetry, and additivity
gives multilinearity.
Intersection with a numerically trivial divisor vanishes after the other
divisors have reduced the cycle to a one-cycle, proving numerical
invariance.
For nonnegativity, first use Chow's lemma to choose a projective
birational morphism \(\mu:Y\to X\) with \(Y\) projective over the ground
field.  The pullbacks of the \(D_i\) are nef, and their intersection is
unchanged by \eqref{eq:appF-degree-formula}.  On the projective variety
\(Y\), approximate the nef classes by ample classes; positivity of ample
intersections and continuity give nonnegativity.
See \cite[Chapters~1 and~19]{Ful}.
\end{proof}

\begin{warning}{Weil is not Cartier}
\label{warn:appF-weil-not-cartier}
On a normal non-\(\mathbf Q\)-factorial variety, a prime Weil divisor need
not admit a pullback as a Cartier divisor and need not define an element of
\(N^1(X)\).
Notation such as \(f^*D\) or \(D^n\) for an arbitrary Weil divisor is
therefore undefined unless an additional pullback or intersection theory has
been specified.
\end{warning}

\section{The projection formula}
\label{sec:appF-projection}

\begin{theorem}\label{thm:appF-projection}
Let \(f:Y\to X\) be proper, let \(D_1,\ldots,D_r\) be Cartier divisors on
\(X\), and let \(\gamma\in A_k(Y)_{\mathbf R}\).
Then
\[
 f_*\bigl(
 f^*D_1\cdots f^*D_r\cap\gamma
 \bigr)
 =
 D_1\cdots D_r\cap f_*\gamma.
\tag{F.3}\label{eq:appF-projection}
\]
Here \(f^*D_i\) acts as \(c_1(f^*\mathcal O_X(D_i))\).
This operator is defined even when the image of \(f\) lies in the
support of \(D_i\), so the local equations do not pull back to a
Cartier divisor.
\end{theorem}

\begin{proof}
For one line bundle \(L\), the first-Chern-class projection formula is
\cite[Proposition~2.3(c)]{Ful}; see also
\cite[Tag~02SU]{StacksProject}:
\[
 f_*\bigl(c_1(f^*L)\cap\gamma\bigr)
 =c_1(L)\cap f_*\gamma.
\]
This is the foundational cycle-level input.  Apply it successively with
\(L=\mathcal O_X(D_i)\), starting with the remaining first-Chern
operators applied to \(\gamma\).  Each application removes one pulled
back operator and gives \eqref{eq:appF-projection} after \(r\) steps.
No pullback of a vanishing local equation is used.
\end{proof}

\begin{corollary}\label{cor:appF-generically-finite}
Let \(f:Y\to X\) be a proper generically finite morphism of degree \(d\)
between proper integral \(n\)-folds.
For \(\mathbf R\)-Cartier divisors \(D_1,\ldots,D_n\) on \(X\),
\[
                   (f^*D_1\cdots f^*D_n)
                   =d(D_1\cdots D_n).
\tag{F.4}\label{eq:appF-degree-formula}
\]
For a proper birational morphism, \(d=1\).
\end{corollary}

\begin{proof}
One has \(f_*[Y]=d[X]\).
Apply \eqref{eq:appF-projection} and take degrees of the resulting
zero-cycles.
\end{proof}

\begin{warning}{Domain of the formula}
\label{warn:appF-projection-domain}
Properness is used to push cycles and take degrees.
Generic finiteness is used only for \(f_*[Y]=d[X]\).
Formula \eqref{eq:appF-degree-formula} does not say that an arbitrary
divisor on \(Y\) is a pullback, and it does not define pullback of a
non-Cartier Weil divisor on \(X\).
\end{warning}

\begin{example}[Blowup test]
\label{ex:appF-blowup}
Let \(\pi:Y\to\mathbf P^2_k\) be the blowup of a \(k\)-rational point
over the ground field \(k\), let
\(H=\pi^*c_1(\mathcal O(1))\), and let \(E\) be the exceptional curve.
Then
\[
                         H^2=1,\qquad H\cdot E=0,\qquad E^2=-1.
\tag{F.5}\label{eq:appF-blowup-table}
\]
The first equality is \eqref{eq:appF-degree-formula}.
The second holds because \(\pi^*\mathcal O(1)|_E\) is trivial.
For the third, \(\mathcal O_E(E)\simeq\mathcal O_{\mathbf P^1}(-1)\).
This table detects a false assertion that every effective divisor has
nonnegative self-intersection.
\end{example}

\section{The birational model category}
\label{sec:appF-models}

Assume now that \(X\) is normal and projective.
A model over \(X\) is a proper birational morphism
\(\pi_Y:Y\to X\) with \(Y\) normal and projective.
We write \(Y'\succeq Y\) when there is a proper birational morphism
\(\mu:Y'\to Y\) over \(X\).

\begin{lemma}\label{lem:appF-common-domination}
Any two models \(Y_1,Y_2\) admit a common normal projective domination.
\end{lemma}

\begin{proof}
Take the closure of the graph of the birational map
\(Y_1\dashrightarrow Y_2\) inside
\(Y_1\times_XY_2\), choose the irreducible component dominating \(X\), and
normalize it.
The two projections are proper and birational.  The graph closure is
projective over \(X\), and its normalization is finite because it is of
finite type over a field; hence the normalized graph is again projective
over \(X\).
\end{proof}

Thus the model category is cofiltered.
No resolution of singularities is used in
\Cref{lem:appF-common-domination}.

\section{Cartier and Weil b-divisors}
\label{sec:appF-bdivisors}

\begin{definition}[Weil b-divisor]
\label{def:appF-weil-bdivisor}
A real Weil b-divisor over \(X\) is a family
\[
                B=(B_Y)_Y,\qquad B_Y\in\Div(Y)_{\mathbf R},
\]
such that for every domination \(\mu:Y'\to Y\),
\[
                              \mu_*B_{Y'}=B_Y.
\tag{F.6}\label{eq:appF-weil-compatibility}
\]
The divisor \(B_Y\) is the trace of \(B\) on \(Y\).
\end{definition}

\begin{definition}[Cartier b-divisor]
\label{def:appF-cartier-bdivisor}
A Cartier b-divisor is a Weil b-divisor \(C\) for which there is a model
\(Y\) and an \(\mathbf R\)-Cartier divisor \(C_Y\) such that
\[
                             C_{Y'}=\mu^*C_Y
\tag{F.7}\label{eq:appF-cartier-determination}
\]
on every \(Y'\succeq Y\).
The pair \((Y,C_Y)\) is a determination of \(C\).
\end{definition}

On a model not dominating a chosen determination, take a common domination
and push down.
The projection formula shows that the result is independent of the common
domination.
Consequently
\[
\begin{aligned}
 \operatorname{CDiv}_{\mathbf R}(\mathfrak X)
   &=\varinjlim_Y\operatorname{CDiv}_{\mathbf R}(Y)
     &&\text{under pullback},\\
 \operatorname{Div}_{\mathbf R}(\mathfrak X)
   &=\varprojlim_Y\operatorname{Div}_{\mathbf R}(Y)
     &&\text{under pushforward}.
\end{aligned}
\tag{F.8}\label{eq:appF-direct-inverse}
\]

\begin{proposition}\label{prop:appF-coefficients}
A Weil b-divisor is equivalent to an assignment \(E\mapsto b_E\in\mathbf R\)
to every prime divisor over \(X\), subject to the condition that on each
model only finitely many divisors occurring on that model have nonzero
coefficient.
Its trace is
\[
                              B_Y=\sum_{E\subset Y}b_EE.
\tag{F.9}\label{eq:appF-trace-coefficients}
\]
\end{proposition}

\begin{proof}
Pushforward preserves the coefficient of a prime divisor that survives and
kills exceptional divisors.
Thus a compatible family gives a model-independent coefficient.
Conversely, the locally finite trace condition makes
\eqref{eq:appF-trace-coefficients} a Weil divisor on every model, and the
traces are compatible under pushforward.
\end{proof}

\begin{example}[A non-Cartier Weil b-divisor]
\label{ex:appF-non-cartier-bdivisor}
On a smooth projective surface, choose infinitely many distinct closed
points \(p_j\).  Let \(Y_j\) be the blowup of the first \(j\) points:
\[
 \cdots\longrightarrow Y_3\longrightarrow Y_2\longrightarrow Y_1
 \longrightarrow X
\]
let \(E_j\) be the prime over \(p_j\), identified with its strict
transform on every higher model where it occurs.
Assign coefficient one to every \(E_j\) and zero to every other prime over
\(X\).
Every normal surface model has only finitely many exceptional primes,
so each fixed trace is finite and the traces push forward compatibly.
For any proposed determination \(Z\to X\), the map is an isomorphism
outside a finite set of closed points.  Choose \(p_j\) outside that set.
The trace \(B_Z\) vanishes near \(p_j\), since every nonzero assigned
coefficient lies on a prime exceptional over \(X\).  Blowing up \(Z\)
at this point therefore gives coefficient zero on the pullback of
\(B_Z\), whereas the assigned coefficient of \(E_j\) is one.
Thus no finite model determines this Weil b-divisor.
\end{example}

\begin{warning}{Two orders}
\label{warn:appF-two-orders}
Coefficientwise effectivity of \(B-C\) is an order on Weil b-divisors.
Pseudo-effectivity of numerical traces is an order on numerical b-divisor
classes.
They are not definitions of one another.
The negativity lemma can relate them for exceptional divisors under stated
hypotheses, but no such relation is formal.
\end{warning}

\section{Intersection products on the Cartier locus}
\label{sec:appF-bdivisor-intersection}

\begin{theorem}\label{thm:appF-cartier-bdivisor-intersection}
Let \(C_1,\ldots,C_n\) be Cartier b-divisors.
Choose a common model \(Y\) determining all of them and set
\[
                         (C_1\cdots C_n)_{\mathfrak X}
                         :=(C_{1,Y}\cdots C_{n,Y}).
\tag{F.10}\label{eq:appF-cartier-bdivisor-product}
\]
This number is independent of \(Y\), symmetric, and multilinear.
\end{theorem}

\begin{proof}
Common determinations exist by
\Cref{lem:appF-common-domination}.
If \(\mu:Y'\to Y\) is a higher model, then every trace on \(Y'\) is the
pullback of its trace on \(Y\).
The birational case of
\eqref{eq:appF-degree-formula} gives equality of the intersections.
Symmetry and multilinearity follow from
\Cref{prop:appF-mixed-cartier}.
\end{proof}

\begin{proposition}\label{prop:appF-one-weil-factor}
Let \(B\) be a Weil b-divisor and let
\(C_1,\ldots,C_{n-1}\) be Cartier b-divisors.
On a common determination \(Y\) of the \(C_i\), define
\[
 (B\cdot C_1\cdots C_{n-1})_{\mathfrak X}
 :=
 \deg\bigl(
 C_{1,Y}\cdots C_{n-1,Y}\cap[B_Y]
 \bigr).
\tag{F.11}\label{eq:appF-one-weil-product}
\]
This is independent of the chosen determination.
\end{proposition}

\begin{proof}
On \(\mu:Y'\to Y\), compatibility gives
\(\mu_*B_{Y'}=B_Y\), while \(C_{i,Y'}=\mu^*C_{i,Y}\).
Apply \eqref{eq:appF-projection} to the cycle \([B_{Y'}]\) and take degrees.
\end{proof}

\begin{warning}{Two Weil factors}
\label{warn:appF-two-weil-factors}
Formula \eqref{eq:appF-one-weil-product} works because every intersection
operator except the initial cycle is Cartier.
There is no resulting product
\(B_1B_2C_3\cdots C_n\) for arbitrary Weil b-divisors.
On a non-\(\mathbf Q\)-factorial model even the first intersection of the
two Weil traces may be undefined.
\end{warning}

\section{Nef Cartier b-divisors}
\label{sec:appF-nef-cartier}

\begin{definition}[Nef Cartier b-divisor]
\label{def:appF-nef-cartier}
A Cartier b-divisor \(C\) is nef if its trace is nef on one determination.
\end{definition}

\begin{lemma}\label{lem:appF-nef-determination}
If \(\mu:Y'\to Y\) is a projective birational morphism, then
\(C_Y\) is nef if and only if \(\mu^*C_Y\) is nef.
Hence \Cref{def:appF-nef-cartier} is independent of the chosen
determination.
\end{lemma}

\begin{proof}
For an integral curve \(\Gamma\subset Y'\),
\[
 (\mu^*C_Y)\cdot\Gamma=C_Y\cdot\mu_*\Gamma\geq0.
\]
If \(\Gamma\) is contracted, the right side is zero.
Conversely, for every integral curve \(C\subset Y\), choose an integral
curve \(\Gamma\subset Y'\) mapping generically finitely onto \(C\).
Such a curve is obtained by choosing a closed point of the nonempty
generic fibre over \(\eta_C\) and taking its closure; its residue field
is finite over \(k(C)\).
Then
\[
 (\mu^*C_Y)\cdot\Gamma
 =\deg(\Gamma/C)(C_Y\cdot C),
\]
so nefness of the pullback implies nefness of \(C_Y\).  Passing to a
common determination proves the final assertion.
\end{proof}

\begin{corollary}\label{cor:appF-nef-positive}
If \(C_1,\ldots,C_n\) are nef Cartier b-divisors, then
\[
                          (C_1\cdots C_n)_{\mathfrak X}\geq0.
\tag{F.12}\label{eq:appF-nef-positive}
\]
\end{corollary}
\begin{proof}
Choose one projective model \(Y\) determining all the \(C_i\).  By
\Cref{lem:appF-nef-determination}, their traces \(C_{i,Y}\) are nef.
Their mixed intersection is nonnegative by
\Cref{prop:appF-mixed-cartier}.  By
\eqref{eq:appF-cartier-bdivisor-product}, this is precisely
\eqref{eq:appF-nef-positive}.
\end{proof}

\section{Numerical b-divisor classes}
\label{sec:appF-numerical-bclasses}

There are two additional hypotheses behind the numerical inverse-limit
language used in the BJ chapters.
We need a cofinal category on which divisor pushforward stays in the chosen
numerical divisor space, and we need a topology for taking its inverse
limit.

Assume for this section that \(X\) is smooth projective over an
algebraically closed field of characteristic zero, and use smooth projective
models, which are cofinal by resolution.
Set
\[
 N^1_{\mathrm{Cart}}(\mathfrak X)
 :=\varinjlim_YN^1(Y),\qquad
 N^1_{\mathrm{Weil}}(\mathfrak X)
 :=\varprojlim_YN^1(Y),
\tag{F.13}\label{eq:appF-numerical-limits}
\]
with pullback and divisor pushforward respectively.
These maps descend to numerical classes on smooth models, by the Chow
projection formula and numerical equivalence; this is the model category
reviewed in \cite[Sections~1.1--1.2]{DF20}.  No corresponding numerical
pushforward on arbitrary singular models is assumed here.
The inverse limit has the projective-limit topology.
For a numerical b-class \(\alpha\), write \(\alpha\geq0\) when every
trace \(\alpha_Y\) is pseudoeffective; write \(\alpha\leq\beta\) when
\(\beta-\alpha\geq0\).  This is the order used in every decreasing-net
statement below.

\begin{definition}[Nef Weil b-class]
\label{def:appF-nef-weil-class}
A nef Weil b-divisor class is a point of
\(N^1_{\mathrm{Weil}}(\mathfrak X)\) in the closure of the nef Cartier
b-divisor classes.
\end{definition}

This definition uses nets.
It is not equivalent, by formal topology alone, to saying every trace is
nef; pushforward traces can satisfy weaker movable positivity conditions.

\begin{theorem}\label{thm:appF-df-nef-intersection}
In the smooth projective characteristic-zero domain above, nef Weil
b-divisor classes carry a symmetric, multilinear, nonnegative intersection
product
\[
                      (\alpha_1\cdots\alpha_n)_{\mathfrak X}.
\tag{F.14}\label{eq:appF-df-product}
\]
It extends \eqref{eq:appF-cartier-bdivisor-product}, is monotone on the nef
cone, is upper semicontinuous in the projective-limit topology, and is
continuous along decreasing nets for which the order and limit are those of
the nef b-class theory.
\end{theorem}

\begin{proof}
Three foundational results are imported here and are not consequences of
finite-level intersection theory: decreasing nef Cartier approximation
\cite[Theorem~2.1]{DF20}, existence and independence of the nef intersection
\cite[Proposition--Definition~3.1]{DF20}, and upper semicontinuity together
with monotone continuity \cite[Theorem~3.2]{DF20}.  We now derive all the
properties in the statement from those inputs.

For a nef Weil b-class \(\alpha\), let \(\mathcal A(\alpha)\) be the
collection of nef Cartier b-classes \(\theta\) with
\(\theta\geq\alpha\), ordered downwards.  The first foundational input is
the monotone approximation theorem
\cite[Theorem~2.1]{DF20}: \(\mathcal A(\alpha)\) contains a decreasing
net converging to \(\alpha\) in the projective-limit topology.  For
\(\alpha_1,\ldots,\alpha_n\), put
\[
 I(\alpha_1,\ldots,\alpha_n)
 :=
 \inf_{\theta_i\in\mathcal A(\alpha_i)}
       (\theta_1\cdots\theta_n)_{\mathfrak X}.
\]
Every number under the infimum is nonnegative by
\Cref{cor:appF-nef-positive}, so the infimum exists.

We isolate the second foundational input.  The negativity and compactness
argument of \cite[Proposition--Definition~3.1]{DF20} proves that for
arbitrary decreasing Cartier approximations
\(\theta_{i,a}\downarrow\alpha_i\), after passing to the product directed
set,
\[
 I(\alpha_1,\ldots,\alpha_n)
       =\lim_a(\theta_{1,a}\cdots\theta_{n,a})_{\mathfrak X}.
\]
In particular, the limit is independent of every choice.  This is the
point at which finite-level intersection theory alone is insufficient.
Define the left side of \eqref{eq:appF-df-product} to be \(I\).

Symmetry and extension of the Cartier product now follow before taking the
limit.  If \(\alpha_1\leq\beta_1\), then
\(\mathcal A(\beta_1)\subseteq\mathcal A(\alpha_1)\), so the infimum
description gives monotonicity.  Nonnegativity was already noted.
For additivity in the first entry, choose decreasing approximations
\(\theta_a\downarrow\alpha\), \(\eta_a\downarrow\beta\), and
\(\theta_{i,a}\downarrow\alpha_i\) for the remaining entries.  Then
\(\theta_a+\eta_a\downarrow\alpha+\beta\), and finite-level
multilinearity gives
\[
\begin{aligned}
 I(\alpha+\beta,\alpha_2,\ldots,\alpha_n)
 &=\lim_a
   ((\theta_a+\eta_a)\theta_{2,a}\cdots\theta_{n,a})_{\mathfrak X}\\
 &=I(\alpha,\alpha_2,\ldots,\alpha_n)
   +I(\beta,\alpha_2,\ldots,\alpha_n).
\end{aligned}
\]
Positive homogeneity is identical, and symmetry gives multilinearity in
every entry.

If \(\alpha_{i,b}\downarrow\alpha_i\), monotonicity of the infimum formula
gives \(\lim_bI(\alpha_{1,b},\ldots,\alpha_{n,b})\geq
I(\alpha_1,\ldots,\alpha_n)\) at once, since
\(\mathcal A(\alpha_{i,b})\subseteq\mathcal A(\alpha_i)\).  The reverse
inequality does \emph{not} follow from that formula: a nef Cartier
\(\theta\geq\alpha_i\) need not dominate \(\alpha_{i,b}\) for any \(b\),
because \(\alpha_{i,b}-\alpha_i\) is a nonzero pseudoeffective b-class and
\(\theta-\alpha_{i,b}\) may fail to be pseudoeffective on some trace, so
\(\bigcup_b\mathcal A(\alpha_{i,b})\) need not be cofinal in
\(\mathcal A(\alpha_i)\).  Continuity along decreasing nets is therefore
part of what is imported from \cite[Theorem~3.2]{DF20}, alongside upper
semicontinuity.

It remains only upper semicontinuity for unrestricted
projective-limit convergence.  This is the other deep half of
\cite[Theorem~3.2]{DF20}: after normalizing against a fixed ample class,
compactness of nef traces on each model and the negativity estimate give
\[
 \limsup_b I(\alpha_{1,b},\ldots,\alpha_{n,b})
       \leq I(\alpha_1,\ldots,\alpha_n).
\]
Thus all the asserted properties hold.
\end{proof}

\Cref{thm:appF-df-nef-intersection} is not a definition of products
of arbitrary Weil b-divisors and is not asserted over a base where the
required model category and approximation theorem are unavailable.
Here \(X\) has absolute dimension \(n\), and the Dang--Favre product has
exactly \(n\) entries.  In particular, it does not define an absolute
quantity \((\alpha^{n+1})_{\mathfrak X}\).  The \((n+1)\)-fold mixed
intersection numbers used for non-Archimedean model metrics belong instead
to the relative theory over a discrete valuation ring: on an
\((n+1)\)-dimensional model they include an actual vertical divisor factor,
for example \((D\cdot\mathcal L_1\cdots\mathcal L_n)\).  Such a BFJ
pairing is not obtained by adding a formal extra copy of an absolute
b-class to \eqref{eq:appF-df-product}.

\begin{corollary}\label{cor:appF-nef-volume}
For a nef b-class \(\alpha\) in the domain of
\Cref{thm:appF-df-nef-intersection}, its b-divisorial volume is
\[
                         \operatorname{vol}(\alpha)
                         =(\alpha^n)_{\mathfrak X}.
\tag{F.15}\label{eq:appF-nef-volume}
\]
If this number is zero, no positive-volume normalization is permitted.
\end{corollary}
\begin{proof}
Apply \Cref{thm:appF-df-nef-intersection} with
\(\alpha_1=\cdots=\alpha_n=\alpha\).  On a decreasing nef Cartier approximation
\(\alpha_i\to\alpha\), the resulting number is
\((\alpha_i^n)\), the usual volume of the nef class \(\alpha_i\).  Decreasing
continuity in \Cref{thm:appF-df-nef-intersection} identifies its limit with
\((\alpha^n)_{\mathfrak X}\), which is the b-divisorial volume by definition.
Nonnegativity follows from the same theorem.  If the value is zero, division
by it is undefined, so none of the probability normalizations used in the
positive-volume theory applies.
\end{proof}

\section{Coefficient functions and the metric sign}
\label{sec:appF-coefficient-functions}

For a prime divisor \(E\) over \(X\), let \(\ord_E\) be its divisorial
valuation.
The coefficient function of a Weil b-divisor is
\[
                 \psi_B(t\ord_E)=t\,\operatorname{coeff}_E(B),
                 \qquad t>0.
\tag{F.16}\label{eq:appF-coefficient-function}
\]
It is homogeneous and determines \(B\).
For an effective Cartier divisor \(D\) with canonical section \(s_D\),
\[
                         \log|s_D|(v)=-v(D)=-\psi_{\overline D}(v).
\tag{F.17}\label{eq:appF-section-sign}
\]
Thus the section function and the coefficient function have opposite signs.

On a nontrivially valued model, \Cref{app:normalizations} instead
associates to a vertical divisor \(D\) the model function
\[
                              \phi_D=v(D).
\]
Equation \eqref{eq:appF-section-sign} is compatible with that convention:
the model function uses the inverse local generator of
\(\mathcal O(D)\), not the canonical section of \(\mathcal O(D)\).

\section{An ample boundary for numerical detection}
\label{sec:appF-Segre-ample-boundary}
\index{Segre class!big and nef boundary}

The projection formula preserves intersections; it does not make
every polarization detect every exceptional ideal. The example
below is \cite[Example~4.7]{CidRuiz26SegrePublished}. It supplies
the boundary of \Cref{thm:appC-ample-Segre-dependence} within
ordinary smooth projective intersection theory.

\begin{example}
\label{ex:appF-big-nef-Segre}
Work over an algebraically closed field. Blow up a point
$p\in\mathbf P^3$ to obtain $Y$, with exceptional plane
$D\simeq\mathbf P^2$. Choose a line $C\subset D$, and blow up
$C$ to obtain $X$, with second exceptional divisor $E$.
Let $\pi:X\longrightarrow\mathbf P^3$ be the composite and put
$L=\pi^*\mathcal O_{\mathbf P^3}(1)$. This line bundle is big
and nef. For every integer $k\geq2$, the subschemes $E\subset kE$
have different integral-closure ideals, but
\[
 \deg_L s^i(E,X)=\deg_L s^i(kE,X)=0
 \qquad(0\leq i\leq3).
\tag{F.18}\label{eq:appF-big-nef-Segre-zero}
\]
\end{example}

Here the pushforward of the Segre class of an effective Cartier
divisor $A$ on a smooth threefold is
\[
 j_{A*}s(A,X)=(A-A^2+A^3)\cap[X].
\tag{F.19}\label{eq:appF-Cartier-Segre}
\]
The first two relevant degrees are $L^2E$ and $-LE^2$.
Both vanish because $E$ maps to the point $p$; equivalently,
$L|_E$ is trivial. The remaining degree is $E^3$.
Its vanishing does \emph{not} follow from contraction of $E$:
the degree of a zero-cycle is still an ordinary degree.

Put $N=N_{C/Y}$ and let $\eta:E=\mathbf P(N)\to C$ be
the exceptional projective bundle, with the convention in the
cited source. The normal sequence is
\[
 0\longrightarrow\mathcal O_{\mathbf P^1}(1)
 \longrightarrow N
 \longrightarrow\mathcal O_{\mathbf P^1}(-1)
 \longrightarrow0.
\tag{F.20}\label{eq:appF-Segre-normal-sequence}
\]
The first term is $N_{C/D}$; the last is the restriction of
$N_{D/Y}=\mathcal O_D(-1)$. Thus $\deg c_1(N)=0$.
The projective-bundle relation, together with
$\mathcal O_X(E)|_E=\mathcal O_{\mathbf P(N)}(-1)$, gives
$\deg E^3=-\deg c_1(N)=0$.
Only the first Chern class of \eqref{eq:appF-Segre-normal-sequence}
is needed; no splitting is asserted. Replacing $E$ by $kE$
multiplies the three components in
\eqref{eq:appF-Cartier-Segre} by $k,k^2,k^3$, respectively,
and proves \eqref{eq:appF-big-nef-Segre-zero}.

On the normal smooth scheme $X$, the invertible ideals
$\mathcal O_X(-E)$ and $\mathcal O_X(-kE)$ are integrally
closed. Their orders along $E$ are $1$ and $k$, so they remain
distinct after integral closure. This is the information lost
by the displayed $L$-degrees. The example does not question
the projection formula, positivity of nef intersections or the
definition of b-divisor products. It shows why the ample
criterion in \cref{sec:appC-Segre-dependence} cannot simply
be invoked with a big and nef model class.

\subsection{Globally generated Chern classes and volume derivatives}
\label{sec:appF-Chern-volume-bridge}

Global generation supplies positivity for certain algebraic
Chern-cycle polynomials. This is a different question from
pointwise positivity of curvature representatives. The following
is the equal-dimension specialization of Proposition~6.2 in
\cite{CidRuiz26MixedSegreV2}; its global hypotheses are explicit.
Let $Y$ be an irreducible projective variety of dimension $N$
over an algebraically closed field, let $q:Y\to\mathbf P^N$
be proper, and let $\mathcal Q_1,\ldots,\mathcal Q_m$ be
globally generated vector bundles. With
$c_t(\mathcal Q)=\sum_j c_j(\mathcal Q)t^j$, write
\[
\begin{aligned}
 q_*\left(\prod_i c_{t_i}(\mathcal Q_i)\cap[Y]\right)
   &=\sum_{|\alpha|\leq N}
        a_\alpha\mathbf t^\alpha H^{|\alpha|}\cap[\mathbf P^N],\\
 A_N(t_0,\mathbf t)
   &=\sum_{|\alpha|\leq N}
        a_\alpha t_0^{N-|\alpha|}\mathbf t^\alpha.
\end{aligned}
\tag{F.21}\label{eq:appF-Chern-coefficient-polynomial}
\]
The attributed result says that $A_N$ has nonnegative integer
coefficients and is denormalized Lorentzian. This specialization
also avoids negative hyperplane exponents in the coefficient
presentation when the domain has larger dimension than the
projective target.

Here is the geometric mechanism from the source. Choose
surjections $\mathcal O_Y^{\ell_i+1}\to\mathcal Q_i$, where
$\ell_i\geq r_i=\operatorname{rank}\mathcal Q_i$, so their
kernels $\mathcal K_i$ have positive rank. On
$Y\times\prod_i\mathbf P^{\ell_i}$, the irreducible incidence
variety $K=\prod_Y\mathbf P(\mathcal K_i)$ is the regular
zero scheme of the tautological section of
$\bigoplus_i p^*\mathcal Q_i\otimes\mathcal O(e_i)$.
The top-Chern and Whitney formulas express its cycle as
\[
 [K]=\sum_{0\leq\alpha_i\leq r_i}
       \left(\prod_i H_i^{r_i-\alpha_i}
                       p^*c_{\alpha_i}(\mathcal Q_i)\right)
                         \cap[Y\times\prod_i\mathbf P^{\ell_i}].
\tag{F.22}\label{eq:appF-Chern-kernel-incidence}
\]
These regular-section and intersection identities are imported
through the cited source. They do not assert positivity of
individual differential forms.

Put $d=\dim K=N+\sum_i(\ell_i-r_i)$ and send $K$ to
$\mathbf P^N\times\prod_i\mathbf P^{\ell_i}$ by $q$ and
the projective coordinates. If the image $T$ has dimension
$d$, its cycle pushforward is $e[T]$, with $e$ the positive
generic degree. With $H_0,H_i$ the restricted hyperplane
classes on $T$, the source's calculation gives
\[
 \mathcal N(A_N)
  =\frac{e}{d!}\left(\prod_i\partial_{t_i}^{\ell_i-r_i}\right)
         \deg_T\bigl((t_0H_0+\cdots+t_mH_m)^d\bigr).
\tag{F.23}\label{eq:appF-Chern-volume-derivative}
\]
The hyperplane classes are nef, so
\Cref{thm:appD-nef-volume-Lorentzian} and derivative closure
give the stated signature property. If $\dim T<d$, the
proper cycle pushforward is zero and the coefficient polynomial
in \eqref{eq:appF-Chern-coefficient-polynomial} is zero.
That case is handled by the zero-polynomial convention; it
is not called the nonzero degree-$d$ volume of a
smaller-dimensional image.

In \cref{sec:appC-mixed-Segre-zeta}, $Y=B_N$ and $q=b_N$,
with the quotient bundles described there. This produces the
modified numerator's coefficients. It supplies no extra
numerical converse under a merely big and nef polarization,
and it does not define products of arbitrary Weil b-divisors.

\subsection{Total Chern intersections on a projective variety}
\label{sec:appF-total-Chern-volume}
\index{Chern class!volume realization}

The preceding equal-dimension pushforward has a general counterpart
without a map to projective space. We use Theorem~3.1 of
\cite{NguyenDang26RichardsonV2}, over $\mathbf C$ as in the
main theorem. Let $X$ be an irreducible projective $d$-fold and
let $F_1,\ldots,F_n$ be globally generated vector bundles. Define
the top-degree coefficient polynomial by
\[
 T_{X,F}(\mathbf x)
    =\sum_{|\alpha|=d}
       \deg_X\left(\prod_i c_{\alpha_i}(F_i)\cap[X]\right)
                                                     \mathbf x^\alpha.
\tag{F.24}\label{eq:appF-total-Chern-polynomial}
\]
A Chern class outside the rank range is zero. The coefficients
are intersection degrees, not differential-form values.

\begin{theorem}
\label{thm:appF-total-Chern-volume}
With these hypotheses, $\mathcal N(T_{X,F})$ is a realizable
volume polynomial. When it is nonzero, it admits a realization
on a smooth projective variety with semiample divisors.
\end{theorem}

Here is the source's construction and its normalization.
Choose sufficiently many generators in each evaluation sequence
\[
 0\longrightarrow K_i\longrightarrow W_i\otimes\mathcal O_X
            \longrightarrow F_i\longrightarrow0,
 \qquad s_i=\operatorname{rank}K_i\geq1,
 \qquad G_i=K_i^\vee.
\tag{F.25}\label{eq:appF-Chern-evaluation}
\]
The bundles $G_i$ are globally generated. Use the convention
that $\mathbf P(G_i)$ parametrizes rank-one quotients, and put
$Y=\prod_X\mathbf P(G_i)$, $\eta_i=s_i-1$ and
$D=d+|\eta|$. Write $\xi_i=c_1(\mathcal O_{\mathbf P(G_i)}(1))$.
The projective-bundle formula and the evaluation sequence give
\[
 \pi_*(\xi_i^{s_i-1+j})=h_j(G_i)=c_j(F_i),
 \qquad
 c_t(F_i)=c_{-t}(G_i)^{-1}.
\tag{F.26}\label{eq:appF-Chern-projective-push}
\]
The equality is a Chow-class identity with the specified quotient
convention. Expanding the top intersection therefore gives
\[
\begin{aligned}
 \frac1{D!}\deg_Y\bigl((x_1\xi_1+\cdots+x_n\xi_n)^D\bigr)
       &=\mathcal N(\mathbf x^\eta T_{X,F}),\\
 \partial^\eta\mathcal N(\mathbf x^\eta T_{X,F})
       &=\mathcal N(T_{X,F}).
\end{aligned}
\tag{F.27}\label{eq:appF-Chern-shift-removal}
\]
The second equality is termwise differentiation:
$(\alpha_i+\eta_i)!$ is reduced to $\alpha_i!$.
The first line is already a realization because $Y$ is
irreducible projective and the tautological divisors are semiample.
\Cref{thm:appD-covolume-preserver} then removes the shift.
The smooth clause uses characteristic-zero resolution in the
source; the smooth realizing variety after differentiation need
not be this particular $Y$. No positive-characteristic resolution
claim is inferred from the projective-bundle identities.

\subsection{Richardson models and tableau coefficients}
\label{sec:appF-Richardson-volume}
\index{Richardson variety!volume model}
\index{skew Schur polynomial}

Richardson varieties provide explicit instances of
\eqref{eq:appF-total-Chern-polynomial}. For partitions
$\mu\subseteq\lambda$, the ordinary skew Schur polynomial
$s_{\lambda/\mu}(x_1,\ldots,x_n)$ records semistandard tableaux
of shape $\lambda/\mu$, with weakly increasing rows and strictly
increasing columns. Its coefficient of $\mathbf x^\alpha$ is
the skew Kostka number $K_{\lambda/\mu,\alpha}$. For strict
partitions one instead has the usual skew Schur $P$- and
$Q$-functions. Their finite-variable polynomials are related by
\[
 Q_{\lambda/\mu}
    =2^{\ell(\lambda)-\ell(\mu)}P_{\lambda/\mu},
\tag{F.28}\label{eq:appF-skew-P-Q-factor}
\]
where $\ell$ counts nonzero parts. These are symmetric-function
conventions, separate from the modified Pfaffians below.

The following is Theorem~1.1 of \cite{NguyenDang26RichardsonV2}.
\begin{theorem}
\label{thm:appF-skew-Schur-volume}
For every ordinary skew shape, $\mathcal N(s_{\lambda/\mu})$
is realizable volume over $\mathbf C$. For every strict skew
shape, $\mathcal N(P_{\lambda/\mu})$ and
$\mathcal N(Q_{\lambda/\mu})$ are realizable volume over
$\mathbf C$. Every nonzero polynomial has a smooth projective
realization with semiample divisors. Zero specializations are
included, and the empty shape gives the constant one.
\end{theorem}

For the ordinary construction, fit $\lambda$ in an $r\times c$
rectangle and let $S$ and $E$ be the tautological subbundle and
quotient bundle on $\operatorname{Gr}(r,r+c)$. The Schubert
dual index is
$\lambda^\vee=(c-\lambda_r,\ldots,c-\lambda_1)$.
The intersection of $\Omega_{\lambda^\vee}$ with
$\Omega_\mu$ for opposite flags is an irreducible Richardson
variety $R$ of dimension $d=|\lambda|-|\mu|$. Pieri's rule
and Schubert duality give
\[
 s_{\lambda/\mu}(\mathbf x)
   =\sum_{|\alpha|=d}
      \deg_R\left(\prod_i c_{\alpha_i}(E)\cap[R]\right)
                                                  \mathbf x^\alpha.
\tag{F.29}\label{eq:appF-Richardson-Chern}
\]
The quotient is globally generated, so the preceding construction
applies. In this Grassmannian presentation one may use $S^\vee$
in the projective bundles, with shift $r-1$ in each variable.
For example, the two disconnected boxes of $(2,1)/(1)$ give
$s_{(2,1)/(1)}=(x_1+\cdots+x_n)^2$, and its normalization
has square coefficients $1/2$ and mixed coefficients $2$.

For the Lagrangian model, choose $N\geq\max(1,\lambda_1)$.
The tautological quotient has rank $N$, and the corresponding
shift is $N-1$. Its Schubert classes are represented by
\emph{modified} $Q$-Pfaffians: the elementary generator is sent
to $c_j(E)$, not to the ordinary Schur $Q$-function evaluated
literally on the Chern roots. The classical convention and
set-of-parts Poincar\'e complement are recorded in
\cite[\S1 and \S3.1]{KreschTamvakis03LagrangianAuthor}.
The type-C cycle transform produces skew $P$; the factor in
\eqref{eq:appF-skew-P-Q-factor} then gives skew $Q$.

Wang's appendix to \cite{NguyenDang26RichardsonV2} gives a
direct type-D spinor realization of skew $Q$. Here the quotient
has rank $N+1$, the shift is $N$, and its special Chern classes
satisfy $c_j(E_D)=2\tau_j$ for $1\leq j\leq N$, while
$c_{N+1}(E_D)=0$. The integral modified-$P$ convention is
\cite[\S1 and \S3]{KreschTamvakis04OrthogonalPublished}.
Appendix~A.5--A.6 of the preprint extends the cycle and
projective-bundle constructions to algebraically closed fields
of arbitrary characteristic. Those intersection realizations
do not include the characteristic-zero smooth-resolution clause.

In characteristic two, the comparison of the ambient type-C and
type-D models uses finite flat radicial maps. On corresponding
Richardson varieties, put $d=|\lambda|-|\mu|$ and
$b=\ell(\lambda)-\ell(\mu)$; their restricted degrees are
$2^b$ and $2^{d-b}$. The restrictions to singular Richardson
varieties are not asserted flat. Factors of two remain integers
in the Chow groups and rational intersection polynomials;
they are not reduced modulo the characteristic. The regular-section
bridge explains the different shifts, without asserting an
isomorphism of the complex type-C and type-D varieties.

An independent proof of the ordinary Lorentzian conclusion is
\cite{Zhang26SkewSchurV1}. If the $n$-variable skew Schur
polynomial is nonzero and $L\geq\lambda_1$, its complementary
shape $(L^n,\mu)/\lambda$ satisfies
\[
 \mathcal N\left((x_1\cdots x_n)^L
       s_{(L^n,\mu)/\lambda}(x_1^{-1},\ldots,x_n^{-1})\right)
       =\mathcal N(s_{\lambda/\mu}(\mathbf x)).
\tag{F.30}\label{eq:appF-skew-Schur-finite-dual}
\]
Dual Jacobi--Trudi gives the identity; the source realizes the
complementary polynomial as a flagged Schubert specialization
and uses dually Lorentzian preservation. This finite-alphabet
complement is different from $\lambda^\vee$ in the Grassmannian.
Deleting an alphabet variable does not commute with the complement.
The independent route proves Lorentzianity rather than geometric
volume realization, and does not supply the skew $P/Q$ models.

If $s_{\lambda/\mu}\ne0$, its exact support is the set of
lattice points of a permutahedron. Let $c_j$ be the lengths of
the skew columns and put $\kappa_p=\#\{j:c_j\geq p\}$.
The support is the integral polymatroid base set with rank
$A\mapsto\sum_j\min(|A|,c_j)$, and vertex coefficients are
one \cite[\S7]{NguyenDang26RichardsonV2}.
The general skew $P/Q$ support description is implicit in that
source; this column formula is not transferred to it.
Normalized coefficients inherit root-direction log-concavity and
\cref{thm:appD-reverse-KT}. The source's covariance bounds concern
the law proportional to $K_{\lambda/\mu,\alpha}\mathbf z^\alpha/\alpha!$,
not a uniform tableau law, and its two-row Littlewood--Richardson
inequalities concern cumulative multiplicities. None implies
log-concavity of each individual Littlewood--Richardson coefficient
or negative association.

\subsection{Actual quotient flags and Chern flow}
\label{sec:appF-Chern-flow}
\index{Chern flow!generated envelope}
\index{quotient flag!creation state}

The packets of \cref{sec:appD-Chern-packets} require generated
bundles representing products that contain nongenerated kernel lines.
The construction in \cite[\S\S3--4]{NguyenDangWang26ChernFlowV4}
retains actual quotient flags and groups whole rows before using
incidence. A formal list of divisor roots would lose the generation
argument.

For $u\in S_n$ let
$c_a=\#\{b>a:u^{-1}(a)>u^{-1}(b)\}$ and
$B_a=(a,a+1,\ldots,a+c_a-1)$, omitting empty rows.
The word $r(u)=B_{n-1}\cdots B_1=(i_1,\ldots,i_\ell)$ is reduced;
row starts decrease. Begin over
$B=\prod_{j=1}^n\mathbf P^{m_j}\times\mathbf P^{m_z}$
with generated lines $L_j$, an additional generated line $M$,
and the split quotient flag $F_j^{(0)}=L_1\oplus\cdots\oplus L_j$,
with $F_0^{(0)}=0$.
At step $t$, with $i=i_t$, form the rank-two bundle
$K_t=\ker(F_{i+1}^{(t-1)}\to F_{i-1}^{(t-1)})$
and its quotient projective bundle. If $S_t$ is the tautological
kernel line, the new prefix is
$F_i^{(t)}=\pi_t^*F_{i+1}^{(t-1)}/S_t$.
Thus one has the literal exact sequence
\[
 0\longrightarrow S_t\longrightarrow
       \pi_t^*F_{i+1}^{(t-1)}\longrightarrow F_i^{(t)}
       \longrightarrow0.
\tag{F.31}\label{eq:appF-Chern-flag-sequence}
\]
Every prefix remains a quotient of the pulled-back original split
bundle and hence is generated. This does not make $S_t$ or each
successive root generated. The source's row bundles are generated
as whole bundles at row completion; for atoms, the co-row bundles
are generated at row start and use the old upper roots at that time.
The polynomial Gysin divided difference needs no inverse divisor
class $1/(x_i-x_{i+1})$ in the Chow ring.

The identity that collects the kernel factors is
\cite[Proposition~4.1]{NguyenDangWang26ChernFlowV4}.
\begin{proposition}
\label{prop:appF-Chern-flow-envelope}
Let $Y$ be a separated finite-type scheme over a field. A finite
directed multigraph has constant-rank generated bundles $F_v$ at
its vertices and, for each edge $e:v\to w$, an exact sequence
$0\to E_e\to F_w\to F_v\to0$. Choose nonnegative integral edge
multiplicities $\kappa_e$ and root multiplicities $r_v$. If
\[
 b_v=r_v+\sum_{e:t(e)=v}\kappa_e
              -\sum_{e:s(e)=v}\kappa_e\geq0,
\tag{F.32}\label{eq:appF-Chern-divergence}
\]
then $U=\bigoplus_v F_v^{\oplus b_v}$ is generated and
\[
 c_t(U)=\prod_v c_t(F_v)^{r_v}\prod_e c_t(E_e)^{\kappa_e}.
\tag{F.33}\label{eq:appF-Chern-flow-class}
\]
The ranks agree, so the top coefficient is the product of the
top Chern classes of all the bundles on the right, with those
multiplicities. No acyclicity assumption is required.
\end{proposition}

The equality follows in the source by summing actual $K^0$ relations
$[E_e]=[F_{t(e)}]-[F_{s(e)}]$ and applying the total Chern operation.
It is interpreted operationally on Chow homology when necessary.
Equalities of ranks alone do not establish those relations.
All factors in \eqref{eq:appF-Chern-flow-class} use one common
Chern parameter $t$. Independent parameters can destroy positivity:
the evaluation sequence on $\mathbf P^1$ has kernel $\mathcal O(-1)$,
and separate parameters give
$(1+uh)(1-vh)=1+(u-v)h$, whose integral is $u-v$.

The flag graph also retains creation times. For rank $j$ put
$\tau_j(t)=\max(\{0\}\cup\{q\leq t:i_q=j\})$.
Only the passive pullbacks between two creations are identified.
At step $t$, its edge goes from the newly created state
$(i_t,t)$ to the existing state
$(i_t+1,\tau_{i_t+1}(t-1))$; this is precisely
\eqref{eq:appF-Chern-flag-sequence}. Collapsing all states of
one rank can join unrelated bundles.
For the all-kernel canonical word let $m_j$ count occurrences of
colour $j$, $L_j$ be its last time, and $L_j=-\infty$ if absent,
with $L_0=-\infty$. The source's Lemma~4.6 says that terminal
supply $d_j$ balances every state exactly when
\[
 d_j\geq\mathbf1_{\{m_j>0,\ L_j>L_{j-1}\}}\quad(j<n).
\tag{F.34}\label{eq:appF-Chern-terminal-balance}
\]
Rank $n$ has no outgoing edge. Staircase differences suffice.
Supply at the terminal state cannot repair an internal deficit
in a different profile.

For a partition $\mu$ the terminal multiplicities are
$d_j=\mu_j-\mu_{j+1}$, with $\mu_{n+1}=0$.
Let $\epsilon_t\in\mathbf Z_{\geq0}$ be the multiplicity of the
kernel line $S_t$ on the edge created at step $t$.
Under the actual balance condition, Theorem~4.7 constructs the
generated packet bundle by adjoining the whole twisted row bundles
to the divergence bundle. Its rank is
$\ell+|\mu|+\sum_t\epsilon_t$, whereas the tower has relative
dimension $\ell$. Its top-Chern pushforward has the predicted
codimension $|\mu|+\sum_t\epsilon_t$.
Finite complement and cap-aware reciprocity then give
\cref{thm:appD-Chern-packet-volume}.

This last step uses the integral universal incidence bundle,
not an asserted integral general zero scheme. With a generated
bundle $V$ of rank $R$ on an integral projective source, choose
$0\to K\to W\otimes\mathcal O\to V\to0$ with positive kernel
rank. On $\mathbf P(K^\vee)$ the generated tautological class $H$
satisfies $p_*H^{\dim W-1}=c_R(V)$.
Differentiating the normalized divisor-volume polynomial to that
order and setting the corresponding variable to zero extracts
the finite-dual top-Chern insertion without an extra factorial
\cite[Proposition~2.8]{NguyenDangWang26ChernFlowV4}.
A general generated-bundle section can have a reducible zero scheme.
The incidence source remains integral and works over the original
field; the smooth refinement is a characteristic-zero statement.

\subsection{Full Chern moments and positive inverse certificates}
\label{sec:appF-Chern-moments}
\index{Chern moment algebra}\index{inverse Chern certificate}
\index{Hodge--Riemann relations!Chern moments}

The full moment algebra retains base classes that a scalar polynomial
may omit. Its formal construction and positivity are separate steps
\cite[\S6]{NguyenDangWang26ChernFlowV4}.
Let $A$ be a finite graded real Poincar\'e-duality algebra of formal
dimension $d$, with degree on $A^d$, and let $C_{i,a}\in A^a$,
$C_{i,0}=1$, vanish outside $0\leq a\leq d$.
Give each $z_i$ degree one. Define a functional on $A[\mathbf z]$ by
\[
 \Lambda_C(a\mathbf z^\alpha)
     =\deg_A\left(a\prod_i C_{i,\alpha_i}\right)
       \quad\text{if }\deg(a)+|\alpha|=d,
\tag{F.35}\label{eq:appF-Chern-moment-functional}
\]
and zero in other degrees.
\begin{definition}
\label{def:appF-Chern-moment-algebra}
The Chern moment algebra is the marked graded algebra
\[
 R_C=A[\mathbf z]/I_C,\qquad
 I_C=\{p:\Lambda_C(pq)=0\text{ for every }q\in A[\mathbf z]\},
\tag{F.36}\label{eq:appF-Chern-moment-algebra}
\]
with the degree induced by $\Lambda_C$ and the labelled generators
$z_i$. Thus $I_C$ is the multiplication annihilator of the functional.
\end{definition}

This is not its linear kernel, and no radical is taken.
Proposition~6.4 of the source gives perfect complementary pairings,
an injective marked copy of $A$, and uniqueness among duality
algebras with those marked moments. Those assertions are algebraic;
they do not yet assert a Hodge cone.

Put $C_i(t)=\sum_{a=0}^d C_{i,a}t^a$ and choose finite inverse
series $K_i(t)C_i(t)=1$, where
$K_i(t)=1+\sum_{j=1}^{s_i}k_{i,j}t^j$, $s_i\geq1$.
The monic presentation and its fibre-top element are
\[
 B=\frac{A[\xi_1,\ldots,\xi_m]}
 {\bigl(\xi_i^{s_i}+k_{i,1}\xi_i^{s_i-1}
                          +\cdots+k_{i,s_i}\bigr)_i},
 \qquad\eta=\prod_i\xi_i^{s_i-1}.
\tag{F.37}\label{eq:appF-Chern-monic-certificate}
\]
Coefficient extraction of $\eta$ gives the degree in dimension
$d+\sum_i(s_i-1)$ and
$\pi_*\prod_i\xi_i^{s_i-1+\alpha_i}=\prod_i C_{i,\alpha_i}$.
There are no alternating signs or factorials in this extraction.
The source's Proposition~6.6 identifies
$R_C=B/\operatorname{Ann}_B(\eta)$, with weighted degree
$[b]\mapsto\deg_B(\eta b)$.

A positive inverse-Chern certificate additionally requires the
full K\"ahler package on an open convex cone in $B^1$, with each
$\xi_i$ and every specified base direction in its closure.
For a duality algebra $H$ of dimension $D$, that package means
Hard Lefschetz for each interior $L$ and positivity of
$(-1)^j\deg_H(abL^{D-2j})$ on
$\ker(L^{D-2j+1}:H^j\to H^{D-j+1})$, as well as duality.
A monic inverse relation by itself supplies none of these signs.

The descent used here is the classical boundary-polarization
result imported in \cite[Lemma~4.6]{AminiHuhLarson25DecompositionV1}
and \cite[Theorem~3.2]{Cattani08MixedLefschetzV2}.
It replaces a module by its multiplication-annihilator quotient
for a closure class, changes the pairing by insertion of that class,
and lowers the formal dimension by one. In Cattani's grading
the image is regraded as $\widetilde V_l=T V_{l+1}$, with pairing
$\widetilde Q(Tu,Tv)=Q(u,Tv)$; its hypothesis is that
$T+\lambda N_0$ is Lefschetz for every $\lambda>0$.
A closure class plus a positive interior class satisfies that
hypothesis. The underlying Cattani--Kaplan--Schmid descent proof
remains a cited classical input, as recorded in
\cref{sec:appL-Chern-flow-sources}.
Successive descent by the factors of $\eta$ gives the package
on $R_C$ from a positive certificate.

For actual bundles the source constructs that certificate
geometrically. Its Proposition~6.11 and Corollary~6.10 give:
\begin{theorem}
\label{thm:appF-geometric-Chern-moments}
Let $X$ be a smooth integral complex projective $d$-fold and use
$A_X^j=H^{2j}(X,\mathbf R)\cap H^{j,j}(X,\mathbf C)$,
with integration. For globally generated bundles $E_i$, set
$C_{i,a}=c_a(E_i)$. The full moment algebra has the K\"ahler
package for the image of the cone
$\{L+\sum_i t_i z_i:L\in\mathcal K_X,\ t_i>0\}$.
If $r_i=\operatorname{rank}E_i$ and
$\theta=\prod_i c_{r_i}(E_i)\ne0$ in $A_X$, then
$A_X/\operatorname{Ann}_{A_X}(\theta)$, with degree
$[a]\mapsto\int_X\theta a$, has the package in dimension
$d-\sum_i r_i$ for the image of the base K\"ahler cone.
\end{theorem}

The coefficient algebra is real diagonal Hodge cohomology,
not an arbitrary Chow ring assumed to satisfy the standard
conjectures. The proof uses the generated evaluation-kernel duals
of \eqref{eq:appF-Chern-evaluation}, padded to ranks at least two.
On their quotient projective bundles, the tautological
Fubini--Study forms are semipositive and positive in relative
directions. A base K\"ahler form plus positive relative multiples
is positive on every nonzero tangent vector. This constructs the
required open cone and its boundary directions.
The arbitrary-field intersection realization is separately
Proposition~2.10 of the source; it does not use complex cohomology.

For instance $X=\mathbf P^2$, $A_X=\mathbf R[h]/(h^3)$,
$\deg(h^2)=1$, and $E=\mathcal O(1)^{\oplus2}$ give
\[
 R_X(E)=\mathbf R[h,z]/(h^3,\ hz-2h^2,\ z^2-h^2).
\tag{F.38}\label{eq:appF-P2-Chern-moments}
\]
Its Hilbert vector is $(1,2,1)$ and the degree-one pairing is
$\left(\begin{smallmatrix}1&2\\2&1\end{smallmatrix}\right)$.
For $L=h+z$ the primitive class $h-z$ has square $-2$.
The inverse series $1-2ht+3h^2t^2$ reconstructs the moments
algebraically, without asserting evaluation kernels of every
auxiliary rank. In contrast, $\xi^2+h\xi=0$ produces the inverse
array $(1,-h,h^2)$ and scalar polynomial
$x^2/2-xy+y^2/2$; its negative coefficient rules out a
certificate with both displayed directions boundary-polarized.

If a duality algebra $H$ is generated by $H^1$ and chosen classes
$\ell_1,\ldots,\ell_N$ span $H^1$, its degree polynomial
$\deg_H((\sum y_i\ell_i)^D)/D!$ recovers $H$ as its apolar
algebra \cite[Proposition~6.12]{NguyenDangWang26ChernFlowV4}.
Canonical-row towers satisfy the generation condition; ample
integral source directions can complete the displayed packet
directions to recover the full weighted algebra. The spanning
condition cannot be deleted. In particular, the quartic of
\cref{ex:appD-volume-apolar-HR} has a volume model while its
smaller apolar algebra fails a middle Hodge sign.
Fixed nonzero insertion gives an annihilator quotient; general
nonflat coefficient base change is a different operation and
need not commute with the moment construction
\cite[Appendix~D.1]{NguyenDangWang26ChernFlowV4}.

\section{Projection-formula audit}
\label{sec:appF-audit}

\begin{center}
\begin{tabular}{p{.30\textwidth}p{.25\textwidth}p{.34\textwidth}}
\toprule
expression & valid input & reason\\
\midrule
\(f^*D\) &
\(D\) Cartier or \(\mathbf Q\)-Cartier after clearing denominators &
line-bundle class; actual divisor pullback also requires nonvanishing
pulled-back equations\\
\((f^*D_1\cdots f^*D_n)\) &
proper generically finite \(f\), Cartier \(D_i\) &
Chow projection formula\\
\((C_1\cdots C_n)_{\mathfrak X}\) &
Cartier b-divisors &
common determination\\
\((B C_1\cdots C_{n-1})_{\mathfrak X}\) &
one Weil and \(n-1\) Cartier b-divisors &
cycle projection formula\\
\((B_1B_2C_3\cdots C_n)_{\mathfrak X}\) &
not defined in general &
two Weil traces need not intersect\\
\((\alpha_1\cdots\alpha_n)_{\mathfrak X}\) &
nef b-classes in \Cref{thm:appF-df-nef-intersection} &
Dang--Favre approximation\\
\bottomrule
\end{tabular}
\end{center}

\section{Consistency checks}
\label{sec:appF-unit-tests}

\begin{enumerate}
\item Pulling all divisors to a higher birational model must leave
      \eqref{eq:appF-cartier-bdivisor-product} unchanged.
\item A degree-\(d\) finite map multiplies the top intersection of pulled
      back Cartier divisors by \(d\), not by \(d^n\).
\item On the blowup of \(\mathbf P^2\), the table
      \eqref{eq:appF-blowup-table} must be recovered.
\item The b-divisor of
      \Cref{ex:appF-non-cartier-bdivisor} must fail every finite
      determination test.
\item A proposed product of two arbitrary Weil b-divisors must identify a
      theory that defines their intersection; common domination alone is
      insufficient.
\item A volume-zero nef b-class has a valid unnormalized self-intersection
      but no associated probability normalization.
\end{enumerate}

The first four assertions are respectively
\Cref{thm:appF-cartier-bdivisor-intersection,cor:appF-generically-finite,ex:appF-blowup,ex:appF-non-cartier-bdivisor}.
The fifth is the obstruction isolated after
\eqref{eq:appF-one-weil-product}, and the sixth is
\Cref{cor:appF-nef-volume}.  Hence the list records already proved
consequences.

\section*{Source note}
\label{sec:appF-source-note}

The cycle-level projection formula and Cartier intersections follow
\cite{Ful}.
The algebraic b-divisor language and positive products have background in
\cite{BFJ09}; nef b-divisor intersection theory is \cite{DF20}.
The homogeneous BJ realization is cited separately in
\Cref{chap:bj-homogeneity-bdivisors}, because it has additional
smoothness, characteristic, and envelope hypotheses.
